\scanpage{115}
On a donc
\[
P_G=\prod_{x\in\widehat A}T_x
=\mathbb U^{\widehat A}\wedge_{\{\pm1\}^{\widehat A}}
\overbrace{\prod_{a\in\widehat A}\varepsilon_a}^{S\varepsilon_G}
\]
où pour $a\in\widehat A$, $\varepsilon_a$ est l'ens. à deux éléments défini par
\[
\begin{cases}
\varepsilon_a=\widehat{\vec A}_s-\{a\}&\text{si }a\in I\\
\varepsilon_a=\displaystyle\bigwedge_{\widetilde a\text{ sur }a}
\bigl(\widehat{\vec A}_{\unsure{\sigma|o}(\widetilde a)}\setminus\widetilde a\bigr)
&\text{si }a\in A=\widehat A\setminus I
\end{cases}
\]
et on trouve un \deleted{\unsure{applicati}} \inserted{morphisme} canonique
\[
\mu_{\vec A}^{\circ}\times SP_G\longrightarrow SPM_{G^\natural}
\qquad\text{où }\mu_{\vec A}^{\circ}=(]0,1[)^{\vec A}
\]
définissant
\[
\unsure{\Pi_1|\pi_1}(SP_G)=S\mathcal T_G=\mathbb Z^{\widehat A}
\longrightarrow\unsure{\Pi_1|\pi_1}(SPM_{G^\natural})=S\mathcal T_{G^\natural}
\]
et, mieux, un hom. de groupoïdes fondamentaux.

Mais \deleted{$\unsure{P_G}=\unsure{\prod_a\varepsilon_a|\{\prod_a\varepsilon_a\}}$} l'ensemble $\varepsilon_G$
(de cardinal $2^{\operatorname{card} A}$) est contenu dans $SP_G$
\[
S\varepsilon_G\subset SP_G
\]
d'où
\[
\mu_{\vec A}^{\circ}\times S\varepsilon_G\longrightarrow PM_{G^\natural}
\]
et comme $\mu_{\vec A}^{\circ}$ est simplement connexe, on en conclut
\[
S\varepsilon_G\longrightarrow\operatorname{Ob}(S\mathcal T_{G^\natural})
\]
(où $S\mathcal T_{?}$ désigne le groupoïde fondamental de $SP_{\unsure{?}}$)

\deleted{\unsure{J'entends dans ces}}
Désignons par $S_0\mathcal T_G$ le s/groupoïde de $S\mathcal T_G=\unsure{\Pi_1|\pi_1}(P_G\unsure{)}$
induit sur $S\varepsilon_G$, on a donc
\[
S_0\mathcal T_G\longrightarrow S\mathcal T_{G^\natural}.
\]
\mspage{114}
